%! Solving Rational Equations — Unit 5, Lesson 5.6 — Slides (Beamer)
\documentclass[aspectratio=169, 11pt]{beamer}
\usepackage{algebra2-beamer}   % provides \forestheader{} and \sectionlabel[color]{}
% Standard coordinate grid + axes: {xmin}{xmax}{ymin}{ymax}
\newcommand{\lgrid}[4]{%
  \draw[step=1, linegray!45, thin] (#1,#3) grid (#2,#4);
  \draw[->, thick, charcoal] (#1,0)--(#2,0) node[below right, font=\tiny]{$x$};
  \draw[->, thick, charcoal] (0,#3)--(0,#4) node[left, font=\tiny]{$y$};}

\begin{document}

% TITLE SLIDE (CourseName is not defined in beamer, so the name is literal)
{
\setbeamercolor{background canvas}{bg=forest}
\begin{frame}[plain]
  \vspace{2.8cm}\hspace{0.55cm}%
  \begin{minipage}{0.9\textwidth}
    {\color{goldacc}\bfseries\small LESSON 5.6}\\[0.5cm]
    {\color{white}\bfseries\LARGE Solving Rational Equations}\\[0.35cm]
    {\color{white}\bfseries\large and extraneous solutions}\\[1.0cm]
    {\color{forestmid}\small Algebra 2: Shepherd~$\cdot$~Unit 5: Rational Functions}
  \end{minipage}
\end{frame}
}

% HOOK
\begin{frame}[t, plain]
  \forestheader{Same algebra. Different answers.}
  \sectionlabel{Hook}
  \begin{columns}[T]
    \begin{column}{0.47\textwidth}
      \begin{block}{Equation (I)}
        \[ \frac{x^2}{x-2}=\frac{4}{x-2} \]
        \centering Multiply by $x-2$: \; $x^2=4$
      \end{block}
    \end{column}
    \begin{column}{0.47\textwidth}
      \begin{block}{Equation (II)}
        \[ \frac{x^2}{x-3}=\frac{4}{x-3} \]
        \centering Multiply by $x-3$: \; $x^2=4$
      \end{block}
    \end{column}
  \end{columns}
  \vspace{8pt}
  \begin{block}{}
    \centering Both give the candidates $x=2$ and $x=-2$. \\[3pt]
    \textbf{But (I) has one solution and (II) has two.} The algebra cannot tell them apart.
    Only the check can.
  \end{block}
\end{frame}

% THE CHECK TABLE
\begin{frame}[t, plain]
  \forestheader{Put the candidates back in}
  \sectionlabel[goldacc]{And look at row one}
  \centering
  \renewcommand{\arraystretch}{1.5}
  \begin{tabular}{|c|c|c|c|l|}
  \hline
  \textbf{Eq.} & \textbf{$x$} & \textbf{Left} & \textbf{Right} & \textbf{True? False? Neither?} \\ \hline
  (I)  & $2$  & $\frac40$ & $\frac40$ & \textbf{neither --- not a statement} \\ \hline
  (I)  & $-2$ & $-1$ & $-1$ & true \\ \hline
  (II) & $2$  & $-4$ & $-4$ & true \\ \hline
  (II) & $-2$ & $\frac{4}{-5}$ & $\frac{4}{-5}$ & true \\ \hline
  \end{tabular}

  \vspace{8pt}
  \begin{block}{}
    \centering A number that makes the equation \emph{undefined} was never in the running.
    You cannot check a value into an equation that does not exist there.
  \end{block}
\end{frame}

% EXPRESSION VS EQUATION
\begin{frame}[t, plain]
  \forestheader{First: is it an expression or an equation?}
  \sectionlabel{Do not mix these up}
  \begin{columns}[T]
    \begin{column}{0.47\textwidth}
      \begin{block}{Lesson 5.3 --- expression}
        \[ \frac{6}{x}-\frac{2}{x-1} \]
        \textbf{Combine} it. Build the LCD, rewrite each piece over it, keep the denominator.
      \end{block}
    \end{column}
    \begin{column}{0.47\textwidth}
      \begin{block}{Lesson 5.6 --- equation}
        \[ \frac{6}{x}-\frac{2}{x-1}=1 \]
        \textbf{Solve} it. Multiply every term by the LCD and the denominators are gone.
      \end{block}
    \end{column}
  \end{columns}
  \vspace{8pt}
  \begin{block}{}
    \centering\textbf{You may clear denominators only when there is an equals sign.}
    An expression has no sides to multiply.
  \end{block}
\end{frame}

% THE FOUR STEPS
\begin{frame}[t, plain]
  \forestheader{The four steps --- always in this order}
  \sectionlabel[goldacc]{Post this}
  \begin{enumerate}\setlength{\itemsep}{6pt}
    \item \textbf{Factor \& restrict.} Factor every denominator; write the list of excluded values
          \emph{at the top of the page}, before you touch the equation.
    \item \textbf{Multiply by the LCD.} Every term --- including any term with no denominator. You are
          multiplying both \emph{sides}, not just the fractions.
    \item \textbf{Solve} the linear or quadratic equation that is left.
    \item \textbf{Check.} Compare each candidate to the Step 1 list. Discard any match. Verify the
          survivors in the \emph{original}.
  \end{enumerate}
  \vspace{6pt}
  \begin{block}{}
    \centering\textbf{Clearing the denominators solves a \emph{different} equation.
    The check is what brings you back.}
  \end{block}
\end{frame}

% WHY IT HAPPENS
\begin{frame}[t, plain]
  \forestheader{Why extraneous solutions exist}
  \sectionlabel{The Warm-Up, cashed in}
  \begin{columns}[T]
    \begin{column}{0.48\textwidth}
      \begin{block}{Multiply by a nonzero number}
        \[ 5=2 \;\xrightarrow{\;\times 3\;}\; 15=6 \]
        Still false. And you can divide by $3$ to get back --- so nothing changed.
      \end{block}
    \end{column}
    \begin{column}{0.48\textwidth}
      \begin{block}{Multiply by zero}
        \[ 5=2 \;\xrightarrow{\;\times 0\;}\; 0=0 \]
        \textbf{Now it is true.} And you cannot divide by $0$ to get back.
      \end{block}
    \end{column}
  \end{columns}
  \vspace{8pt}
  \begin{block}{}
    The LCD is an \emph{expression}, and at each excluded value it \emph{is} zero. So the cleared
    equation can gain answers the original never had --- and those answers can only be the excluded
    values. \\[3pt]
    \centering\textbf{An extraneous solution is never a random number. It is always on your Step 1
    list.}
  \end{block}
\end{frame}

% WORKED A
\begin{frame}[t, plain]
  \forestheader{Worked Example A: $\dfrac{6}{x}-\dfrac{2}{x-1}=1$}
  \sectionlabel[goldacc]{Nothing gets thrown away}
  \begin{columns}[T]
    \begin{column}{0.48\textwidth}
      \begin{itemize}\setlength{\itemsep}{4pt}
        \item Restrictions: $x\neq0$, $x\neq1$. \; LCD $=x(x-1)$.
        \item Clear: $6(x-1)-2x=x(x-1)$.
        \item $6x-6-2x=x^2-x \Rightarrow x^2-5x+6=0$.
        \item $(x-2)(x-3)=0$, so candidates $2$ and $3$.
      \end{itemize}
    \end{column}
    \begin{column}{0.48\textwidth}
      \begin{block}{Step 4 --- check}
        Neither candidate is on $\{0,1\}$. \\[4pt]
        $x=2$: \; $3-2=1$ \;\checkmark \\[2pt]
        $x=3$: \; $2-1=1$ \;\checkmark \\[4pt]
        \centering Solution set $\{2,3\}$.
      \end{block}
    \end{column}
  \end{columns}
  \vspace{6pt}
  {\footnotesize Notice the term on the right. It has no denominator, and the LCD still multiplies it.
  Forgetting that is the most common wrong answer of the day.}
\end{frame}

% WORKED B
\begin{frame}[t, plain]
  \forestheader{Worked Example B: one survives, one does not}
  \sectionlabel{$\dfrac{x}{x+1}+\dfrac{2}{x-1}=\dfrac{2}{x^2-1}$}
  \begin{columns}[T]
    \begin{column}{0.50\textwidth}
      \begin{itemize}\setlength{\itemsep}{3pt}
        \item $x^2-1=(x-1)(x+1)$, so $x\neq1,-1$.
        \item Clear: $x(x-1)+2(x+1)=2$.
        \item $x^2+x=0 \Rightarrow x(x+1)=0$.
        \item Candidates: $0$ and $-1$.
      \end{itemize}
    \end{column}
    \begin{column}{0.46\textwidth}
      \begin{block}{Step 4}
        $-1$ is on the list $\Rightarrow$ \textbf{extraneous}. \\[3pt]
        $x=0$: \; $0+(-2)=-2$ \;\checkmark \\[3pt]
        \centering Solution set $\{0\}$.
      \end{block}
    \end{column}
  \end{columns}
  \vspace{6pt}
  {\footnotesize You did not have to substitute $-1$ to reject it. Step 4 is reading two short lists and
  looking for overlap.}
\end{frame}

% THE GRAPH
\begin{frame}[t, plain]
  \forestheader{Now see it --- and 5.5 does the work}
  \sectionlabel[goldacc]{Graphical check}
  \begin{columns}[T]
    \begin{column}{0.52\textwidth}
      \centering
      \begin{tikzpicture}[scale=0.30]
        \lgrid{-6}{7}{-4}{6}
        \draw[navy, dashed, thick] (1,-4)--(1,6);
        \draw[navy, dashed, thick] (-6,1)--(7,1);
        \node[navy, font=\tiny, above right] at (1,5.2) {$x=1$};
        \begin{scope} \clip (-6,-4) rectangle (7,6);
          \draw[very thick, forest, domain=-6:0.8, samples=110, smooth] plot (\x,{1+1/((\x)-1)});
          \draw[very thick, forest, domain=1.25:7, samples=110, smooth] plot (\x,{1+1/((\x)-1)});
        \end{scope}
        \draw[forest, very thick, fill=white] (-1,0.5) circle (4.4pt);
        \fill[charcoal] (0,0) circle (3.4pt);
      \end{tikzpicture}
    \end{column}
    \begin{column}{0.46\textwidth}
      \[ \frac{x(x+1)}{(x-1)(x+1)}=\frac{x}{x-1},\; x\neq-1 \]
      \begin{itemize}\setlength{\itemsep}{3pt}
        \item Cancelled factor $\Rightarrow$ \textbf{hole} at $\left(-1,\frac12\right)$.
        \item Surviving factor $\Rightarrow$ wall at $x=1$.
        \item Numerator zero at $x=0$.
      \end{itemize}
    \end{column}
  \end{columns}
  \vspace{6pt}
  \begin{block}{}
    \centering\textbf{True solutions are $x$-intercepts. Extraneous candidates are holes.}
    Yesterday's cancel test, sorting today's answers.
  \end{block}
\end{frame}

% CROSS-MULTIPLY AND NO SOLUTION
\begin{frame}[t, plain]
  \forestheader{Two more things}
  \sectionlabel{Shortcut, and an answer people refuse to write}
  \begin{columns}[T]
    \begin{column}{0.48\textwidth}
      \begin{block}{Cross-multiplying}
        \[ \frac{2}{x-1}=\frac{x}{x+2} \]
        Only for a \emph{proportion}: one fraction $=$ one fraction. \\[3pt]
        $2(x+2)=x(x-1) \Rightarrow x=4$ or $x=-1$. Neither excluded. \\[3pt]
        It is not a new rule --- it is Step 2, pre-cancelled.
      \end{block}
    \end{column}
    \begin{column}{0.48\textwidth}
      \begin{block}{``No solution'' is an answer}
        \[ \frac{1}{x-2}+\frac{1}{x+2}=\frac{4}{x^2-4} \]
        Clears to $2x=4$, so the only candidate is $x=2$ --- which is excluded. \\[3pt]
        Every candidate extraneous $\Rightarrow$ \textbf{no solution}. Write it down and move on.
      \end{block}
    \end{column}
  \end{columns}
\end{frame}

% CLOSE
\begin{frame}[t, plain]
  \forestheader{Today in one line --- and where it goes next}
  \sectionlabel[goldacc]{Close}
  \begin{block}{}
    \centering\textbf{Clearing the denominators solves a different equation --- the check is what
    brings you back.}
  \end{block}
  \vspace{6pt}
  \begin{columns}[T]
    \begin{column}{0.48\textwidth}
      \begin{block}{Errors that cost the most}
        \begin{itemize}\setlength{\itemsep}{2pt}
          \item No restriction list, so no way to check.
          \item Multiplying the fractions but not the whole side.
          \item Cross-multiplying with three terms.
          \item Rejecting an answer for being negative.
          \item Refusing to write ``no solution.''
        \end{itemize}
      \end{block}
    \end{column}
    \begin{column}{0.48\textwidth}
      \begin{block}{Next: Lesson 5.7}
        Direct and inverse variation --- the rational parent $y=\frac kx$ as a model. Rejecting answers
        comes back, but for a new reason: not because the algebra forbids them, because the
        \emph{situation} does.
      \end{block}
    \end{column}
  \end{columns}
\end{frame}

\end{document}
